\section{Introduction}
\label{sec:introduction}

A corridor viewed through an open doorway, the same corridor reflected in a
planar mirror, and a photograph of it displayed on a screen can be arranged to
produce an identical image. The three scenes differ in the surface a hand would
touch, which is empty air in the first case, a pane of glass in the second and a
rigid panel in the third, yet at a single viewpoint the light reaching the sensor
is the same. The physical explanation of the observed geometry is therefore not
determined by the image, and any answer produced from that image alone is
selected by a prior rather than by evidence.

This distinction carries practical weight wherever perception is coupled to
action. A service robot must separate a doorway from its reflection before a path
is committed, an assistive system is expected to report the surface in front of a
user rather than the scene depicted on it, and an augmented-reality device places
content on a panel rather than inside the room that panel shows. In each case the
cost of an error is set by physical contact, so the quantity of interest is the
geometry of the first real surface along the ray, and the failure of interest is a
confident report of geometry that is not present.

\paragraph{From accurate depth to identifiable depth.}
Monocular depth estimation now recovers relative geometry reliably across diverse
scenes, and large pre-trained predictors transfer to unfamiliar imagery with
little adaptation. Several complementary lines of work extend perception to
exactly the configurations described above. Multi-layer and order-aware
formulations represent more than one surface along a ray, so that a transparent
interface and the scene behind it are both retained. Dedicated detectors recover
mirrors and glass as first-class scene elements. Active and next-best-view
perception treats observer motion as a resource to be allocated toward informative
viewpoints. Uncertainty estimation supplies a scalar that accompanies each
prediction, selective prediction converts such a scalar into a decision to answer
or to defer, and split conformal prediction supplies finite-sample guarantees on
how often deferral occurs.

Each of these establishes a capability on which the present work depends, and
together they address what the depth is, which surfaces are present, where the
camera should move next, and how certain a prediction is. The matched
configurations above raise a further question, which is the subject of this paper:
is the physical explanation of this visual geometry recoverable at all from the
observations that this observer is able to make?

\paragraph{Why this is a distinct question.}
Accuracy is a property of a prediction and uncertainty is a property of a
predictor, whereas identifiability is a property of the triple formed by the
competing physical explanations, the actions available to the observer, and the
resolution at which differences become perceptible. A view-tracked display
observed within its aperture is geometrically indistinguishable from the scene it
depicts, and a static planar mirror reflecting a static room produces content
parallax matching that of a real room seen through an opening of the same shape.
In both cases the answer is not merely difficult to obtain, but absent from the
available evidence. Enlarging a model or collecting further images from the same
viewpoint leaves the situation unchanged, whereas a suitable change of viewpoint
resolves it immediately. Identifiability is defined here relative to an action
set for this reason, and the label ``not identifiable'' is treated as a statement
about that triple rather than as a further physical hypothesis. The framing
yields a prediction that can be tested independently of architecture: where the
evidence cannot decide, a scalar confidence has no mechanism by which to register
the absence of evidence, because the same image remains consistent with every
hypothesis.

\paragraph{Contributions.}
\begin{enumerate}\itemsep2pt
\item \textbf{Action-set identifiability as a decision rule.} Separability between
two image-formation hypotheses is defined as the distance between the consequences
they predict for an intervention, and identifiability as the best such separability
attainable over an available action set. Falling below a perceptual threshold
licenses abstention, which turns recoverability into an explicit and computable
precondition for answering.

\item \textbf{A matched-counterfactual benchmark whose premise is measured rather
than assumed.} Competing mechanisms are pixel-identical at the reference view to a
maximum deviation of order $10^{-14}$~px, so chance-level single-view performance
is a verified property of the construction. Cases that no permitted action
resolves are included by design, which makes correct abstention a scorable outcome.

\item \textbf{A weakest-pair criterion for selecting the intervention.} Actions
are chosen to maximise the smallest separability among hypothesis pairs still in
contention, rather than a belief-weighted sum over all pairs. The criterion is
$T$-optimality from optimal design for model discrimination
\citep{atkinson1975design,hunter1965experimental}, applied here to active visual
perception and compared against the summed objective on this problem class.

\item \textbf{A protocol for real photographs, and the ground-truth admissibility
criterion it requires.} Calibrated stereo is treated as a one-element action set,
predictions are calibrated outside the ambiguous region and scored inside it, and
images whose reference measurement cannot arbitrate the question are identified
and excluded.

\item \textbf{Evidence across model families, and a decodable signal.} Behaviour
is characterised over published checkpoints spanning several architecture families,
with trivial baselines reported alongside the proposed measures, and a policy
fitted on features pooled from a frozen encoder transfers to held-out
architectures.
\end{enumerate}
